\section{¿Qué pasaría si...? Ideas especulativas para el proyecto}

\subsection{¿Qué pasaría si Apophis tuviera masa comparable a la de la Luna?}

Simular el encuentro con una masa artificial mucho mayor para Apophis, del orden de la masa lunar. Estudiar a partir de qué escala de masa la gravedad del asteroide empieza a perturbar apreciablemente la órbita de la Tierra durante el encuentro. Evaluar además si existe una masa crítica para la cual Apophis podría quedar capturado temporal o permanentemente por el sistema Tierra--Sol.

\subsection{¿Qué pasaría si el encuentro ocurriera exactamente en el plano de la eclíptica?}

Forzar $I = 0^{\circ}$ en los elementos orbitales de Apophis y repetir la simulación. Analizar si la distancia mínima de aproximación cambia de manera significativa y explorar si existe alguna inclinación crítica que maximice o minimice la distancia de encuentro.

\subsection{¿Qué pasaría si Júpiter no existiera?}

Re-simular los últimos 100 años de la órbita de Apophis eliminando a Júpiter del sistema. Comparar el estado orbital resultante con el escenario nominal y determinar si el encuentro de 2029 habría ocurrido con una geometría distinta. Esta prueba permite cuantificar cuánto influye Júpiter como perturbador dominante del sistema solar interior en este caso particular.

\subsection{¿Qué pasaría si el encuentro hubiera ocurrido del lado nocturno de la Tierra?}

Tener en cuenta la rotación terrestre durante el paso cercano y calcular qué punto geográfico de la Tierra quedaría bajo Apophis en el instante de máxima aproximación. A partir de ello, identificar desde qué regiones o ciudades del mundo el asteroide sería visible a simple vista y comparar ese escenario con la geometría real del encuentro.

\subsection{¿Qué pasaría si Apophis pasara exactamente a la distancia de los satélites geoestacionarios?}

Ajustar artificialmente la órbita de Apophis para que la distancia mínima a la Tierra sea exactamente $42{,}164\,\mathrm{km}$. Calcular la deflexión angular asociada a ese paso cercano y el cambio resultante en su período orbital heliocéntrico. Como extensión, evaluar si después de esa perturbación seguiría clasificándose como un asteroide tipo Aten.

\subsection{¿Qué pasaría si el encuentro ocurriera 6 horas antes o 6 horas después?}

Desplazar el estado inicial de Apophis para forzar dos encuentros alternativos: uno 6 horas antes y otro 6 horas después del escenario nominal. Comparar distancia mínima, velocidad relativa y ubicación del subpunto sobre la Tierra en los tres casos. Así se puede mostrar cuánto cambia la geografía de visibilidad por desplazamientos temporales moderados y físicamente plausibles.

\subsection{¿Qué pasaría si la Tierra tuviera una masa un 5\% mayor?}

Modificar únicamente la masa terrestre en la simulación y repetir el paso cercano. Cuantificar cómo cambian la deflexión gravitacional, el tiempo dentro de la esfera de influencia y los elementos orbitales posteriores de Apophis. Es un experimento hipotético claro para estudiar sensibilidad dinámica frente al parámetro gravitacional del planeta.

\subsection{¿Qué pasaría si Apophis rotara muy rápido y expulsara material superficial?}

Suponer un escenario de actividad débil inducida por rotación rápida (eyección de polvo o fragmentos pequeños) durante el encuentro. Modelar esas partículas como trazadores con velocidades de eyección bajas respecto al cuerpo principal y seguir su evolución por algunos días. Evaluar si alguna nube de partículas produciría una firma observacional distinta del núcleo principal.

\subsection{¿Qué pasaría si se aplicara una maniobra cinética mínima décadas antes?}

Introducir una perturbación impulsiva muy pequeña (orden de mm/s a cm/s) en una fecha anterior al 2029 y estudiar cuánto cambia la distancia mínima del encuentro. Repetir para varios años de aplicación del impulso y construir una curva ``anticipación vs. eficacia'' que muestre por qué correcciones minúsculas, si se hacen temprano, producen diferencias macroscópicas en la geometría final.

\subsection{¿Qué pasaría si el encuentro se analizara en un modelo Tierra--Luna explícito de alta resolución temporal?}

Repetir la aproximación con pasos de integración muy finos e incluyendo de forma explícita la Luna como tercer cuerpo cercano dominante en la fase crítica. Comparar la distancia mínima y la velocidad relativa con un modelo simplificado sin Luna en esa escala temporal. El objetivo es detectar si aparecen micro-perturbaciones medibles en la trayectoria nominal durante las horas más cercanas.
